\documentclass[12pt, letterpaper]{article}

% ── Encoding & fonts ──────────────────────────────────────────────────────────
\usepackage[T1]{fontenc}
\usepackage[utf8]{inputenc}

% ── Page layout ───────────────────────────────────────────────────────────────
\usepackage[margin=1in]{geometry}

% ── Math ──────────────────────────────────────────────────────────────────────
\usepackage{amsmath, amssymb}

% ── Tables ────────────────────────────────────────────────────────────────────
\usepackage{booktabs}
\usepackage{array}

% ── Figures & captions ────────────────────────────────────────────────────────
\usepackage{graphicx}
\usepackage{caption}
\captionsetup{font=normalsize, labelfont=bf, labelsep=period}

% ── Lists ─────────────────────────────────────────────────────────────────────
\usepackage{enumitem}
\setlist[itemize]{noitemsep, topsep=4pt, leftmargin=1.5em}

% ── Code / monospace ──────────────────────────────────────────────────────────
\usepackage{listings}
\lstset{basicstyle=\ttfamily\normalsize, breaklines=true}

% ── Header / footer ───────────────────────────────────────────────────────────
\usepackage{fancyhdr}
\pagestyle{fancy}
\fancyhf{}
\fancyhead[L]{\normalsize Opinion Mining in Amazon Reviews}
\fancyhead[R]{\normalsize Machine Learning 2026}
\fancyfoot[C]{\normalsize\thepage}
\renewcommand{\headrulewidth}{0pt}
\renewcommand{\footrulewidth}{0pt}

% ── Hyperlinks — black, no colored boxes ──────────────────────────────────────
\usepackage[colorlinks=false, pdfborder={0 0 0}]{hyperref}

% ── Paragraph spacing ─────────────────────────────────────────────────────────
\usepackage{parskip}
\setlength{\parskip}{6pt}

% ── Section formatting — black, bold, no decorators ──────────────────────────
\usepackage{titlesec}
\titleformat{\section}{\large\bfseries}{\thesection.}{0.5em}{}
\titleformat{\subsection}{\normalsize\bfseries}{\thesubsection}{0.5em}{}
\titleformat{\subsubsection}{\normalsize\itshape}{\thesubsubsection}{0.5em}{}
\titlespacing{\section}{0pt}{14pt}{6pt}
\titlespacing{\subsection}{0pt}{10pt}{4pt}

% ── Table column helpers ──────────────────────────────────────────────────────
\newcolumntype{C}[1]{>{\centering\arraybackslash}p{#1}}
\newcolumntype{L}[1]{>{\raggedright\arraybackslash}p{#1}}

% ─────────────────────────────────────────────────────────────────────────────
\begin{document}

% ── Title page ────────────────────────────────────────────────────────────────
\begin{titlepage}
  \centering
  \vspace*{2cm}
  {\large\bfseries Universidad Autónoma de Tamaulipas\\[4pt]}
  {\normalsize Facultad de Ingeniería y Ciencias\\[2.5cm]}

  {\LARGE\bfseries Opinion Mining in Amazon Product Reviews\\[8pt]}
  {\large Sentiment Classification with TF-IDF and Word2Vec
    (CBOW\,/\,Skip-gram)\\[1.8cm]}

  {\normalsize\itshape Machine Learning Project\\[1.5cm]}

  {\large\bfseries Hugo Emiliano Vargas Briones\\[6pt]}
  {\normalsize Machine Learning --- Dr.\ José Lázaro Martínez Rodríguez\\[6pt]}
  {\normalsize April 28, 2026}

  \vfill
\end{titlepage}

% ── Table of contents ─────────────────────────────────────────────────────────
\tableofcontents
\newpage

% ─────────────────────────────────────────────────────────────────────────────
\section{Introduction}

Amazon generates millions of product reviews every year. Most of that content
goes unused: the star rating reaches recommendation algorithms, but the text
below---where users explain in their own words what worked, what failed, which
brand they bought and whether they would return---is rarely processed
systematically. Reading 58,000 reviews by hand is not a realistic option.

This project applies text-mining and machine-learning techniques to the
\textbf{Seeds \& Grasses} subcategory of the Amazon Reviews 2023 dataset,
comprising 58,108 English-language reviews. The objective is to build an
automatic five-category sentiment classifier and compare the performance of
three different text representations: TF-IDF in three $n$-gram configurations
and Word2Vec in two architectural variants, CBOW and Skip-gram.

The original project instructions specified the concrete-spalling dataset;
however, at Dr.\ Lázaro's indication the Amazon Reviews 2023 corpus was used
instead, since it connects directly to an ongoing thesis project on sentiment
analysis.

% ─────────────────────────────────────────────────────────────────────────────
\section{Theoretical Background}

\subsection{Text Preprocessing}

Before vectorization, each review undergoes a transformation pipeline:
noise removal (HTML tags, URLs, punctuation), tokenization, lemmatization, and
stop-word elimination. Lemmatization reduces every word to its canonical form,
unifying inflected variants of the same term~\cite{spacy}. No minimum-length
filter is applied to tokens; short particles such as negations can be
decisive for polarity and should be preserved.

\subsection{Vector Representations}

\paragraph{TF-IDF.}
Term Frequency--Inverse Document Frequency weights each term by its frequency
in the document and its rarity across the corpus~\cite{sklearn}. Three
$n$-gram configurations were evaluated: unigrams only \mbox{(1,\,1)},
unigrams and bigrams \mbox{(1,\,2)}, and unigrams, bigrams, and trigrams
\mbox{(1,\,3)}. Trigrams capture multi-word expressions such as
\textit{``very well done''} or \textit{``not good at all''} that unigrams
alone would miss.

\paragraph{Word2Vec.}
A dense vector representation that places semantically related words close
together in a continuous embedding space~\cite{word2vec}. Two architectures
were evaluated:

\begin{itemize}
  \item \textbf{CBOW} (\texttt{sg=0}): predicts the central word from its
    surrounding context. Learns well-represented frequent vocabulary efficiently.
  \item \textbf{Skip-gram} (\texttt{sg=1}): predicts the context from the
    central word. Tends to produce better representations for infrequent words,
    which is relevant given the technical agronomic vocabulary of this corpus.
\end{itemize}

Reviews are represented as the average of their constituent word vectors.
Both models are trained exclusively on training-set tokens to prevent data
leakage. Vectors are scaled to $[0,\,1]$ with a \texttt{MinMaxScaler}
fitted on the training set only.

\subsection{Classification Algorithms}

Six supervised classifiers were evaluated:

\begin{itemize}
  \item \textbf{Naïve Bayes}: probabilistic classifier assuming feature
    independence. Fast and effective on text. Key hyperparameter:
    \texttt{var\_smoothing}.
  \item \textbf{Logistic Regression}: linear model estimating class
    probabilities. Performs well in high-dimensional spaces~\cite{lovera2023}.
    Key hyperparameter: $C$ (inverse L2 regularization strength).
  \item \textbf{$k$-NN}: classifies by the majority class among the $k$
    nearest neighbours. Key hyperparameter: $k$.
  \item \textbf{Decision Tree}: builds binary classification rules
    recursively~\cite{sklearn}. Key hyperparameters: \texttt{max\_depth}
    and \texttt{criterion} (Gini or entropy).
  \item \textbf{SVM}: finds the maximum-margin hyperplane separating classes;
    with an RBF kernel it handles non-linearly separable data~\cite{cortes1995}.
    Key hyperparameters: $C$ and \texttt{kernel}.
  \item \textbf{MLP}: dense neural network with hidden layers and ReLU
    activations~\cite{sklearn}. Key hyperparameters: layer architecture and
    L2 regularization (\texttt{alpha}).
\end{itemize}

\subsection{Evaluation Metrics}

Because the class distribution is imbalanced, all metrics are
\textbf{macro-averaged}: each class contributes equally regardless of its
frequency.

\begin{itemize}
  \item \textbf{Accuracy}: fraction of correct predictions.
  \item \textbf{Precision}: of all samples predicted as class $c$, how many
    actually belong to $c$.
  \item \textbf{Recall}: of all true members of class $c$, how many are
    correctly predicted.
  \item \textbf{F1 macro}: harmonic mean of precision and recall, averaged
    across classes.
  \item \textbf{CV F1}: mean macro F1 from $k$-fold cross-validation on the
    training set. This is the \emph{primary} generalisation indicator; the
    train--test gap is secondary when CV F1 is stable across folds.
\end{itemize}

% ─────────────────────────────────────────────────────────────────────────────
\section{Methodology}

\subsection{Corpus and Labelling}

The Seeds \& Grasses subcategory of the Amazon Reviews 2023 dataset was used.
Unlike the previous version of this project---which derived labels directly
from star ratings---sentiment categories were assigned by an \textbf{ensemble
of three local language models}: \texttt{qwen2.5:1.5b}, \texttt{gemma2:2b},
and \texttt{phi4-mini}, all running locally via Ollama. Each model labelled
every review independently; the final label is determined by majority vote.
Records where at least two of the three models agreed (concordance $\geq 0.67$)
were retained, yielding 54,084 records with a guaranteed minimum label quality.

The five sentiment categories are:
\textit{very bad}, \textit{bad}, \textit{neutral}, \textit{good}, and
\textit{very good}.
This five-level scheme captures the distribution of real user experiences more
accurately than a binary or three-class star-based scheme.

\subsection{Stratified Sample}

A \textbf{stratified sample of 6,000 reviews} (1,200 per class) was used for
the demonstration run on a MacBook Air M4. The train/test split (80\,\%/20\,\%)
is performed with stratification, guaranteeing the same class proportions in
both subsets. The full run (${\sim}54{,}000$ records) executes on the research
server (Intel Xeon Silver 4210R, 20 cores, RTX 4060, 48\,GB RAM).

\subsection{Preprocessing}

Each review is transformed through: (1)~lower-casing and removal of HTML
tags, URLs, and non-alphabetic characters; (2)~morphological analysis with
spaCy (\texttt{en\_core\_web\_sm}) using \texttt{nlp.pipe} in batches of 512
records for efficiency; (3)~stop-word removal. No minimum-length token filter
is applied, preserving semantically relevant short particles. Named-entity
recognition and topic modelling are not part of this experiment.

\subsection{Split and Vectorisation}

\textbf{The correct order is: split first, then vectorise.} If the vectoriser
is fitted before the split, the test set sees its own statistics and evaluation
metrics become inflated. TF-IDF and Word2Vec are fitted exclusively on
training data; the test set is only transformed.

Three representations are compared:

\begin{itemize}
  \item \textbf{TF-IDF (1,3)}: unigrams, bigrams, and trigrams.
    \texttt{max\_features=3,000} (demo); \texttt{max\_df=0.90};
    \texttt{min\_df=3}.
  \item \textbf{Word2Vec CBOW} (\texttt{sg=0}): predicts the central word
    from context. \texttt{vector\_size=100}, \texttt{window=5},
    \texttt{min\_count=2}, \texttt{epochs=10}, \texttt{workers=4}.
  \item \textbf{Word2Vec Skip-gram} (\texttt{sg=1}): predicts context from
    the central word. Same hyperparameters as CBOW.
\end{itemize}

\subsection{Hyperparameter Search}

A \texttt{GridSearchCV} with \textbf{3 stratified folds} over the training
set (demo configuration) was used to select hyperparameters. The selection
criterion was macro F1. Models supporting \texttt{class\_weight} received the
value \texttt{'balanced'} as a fixed parameter. The MLP receives per-sample
weights computed with \texttt{compute\_sample\_weight}. The best
hyperparameters found with TF-IDF are reused for all three vectorisations.

% ─────────────────────────────────────────────────────────────────────────────
\section{Experiments and Results}

\subsection{Experimental Configuration}

\begin{table}[h!]
\centering
\caption{General experimental parameters.}
\label{tab:config}
\begin{tabular}{L{4.5cm} L{9cm}}
\toprule
\textbf{Parameter} & \textbf{Value} \\
\midrule
Corpus & Seeds \& Grasses (Amazon Reviews 2023) \\
Sample size & 6,000 reviews (demo) $\cdot$ 54,000 full run \\
Classes & very bad $\cdot$ bad $\cdot$ neutral $\cdot$ good $\cdot$ very good \\
Labelling & LLM ensemble (qwen2.5:1.5b + gemma2:2b + phi4-mini) \\
Split & 80\,\% training (4,800) $\cdot$ 20\,\% test (1,200) \\
Cross-validation & StratifiedKFold $k=5$ (demo) $\cdot$ $k=10$ full run \\
TF-IDF & 3 configurations: (1,\!1) $\cdot$ (1,\!2) $\cdot$ (1,\!3) \\
Word2Vec & CBOW (\texttt{sg=0}) and Skip-gram (\texttt{sg=1}), 100 dims \\
\bottomrule
\end{tabular}
\end{table}

\subsection{Hyperparameter Search Results}

Table~\ref{tab:grid} summarises the explored parameters and the best
configuration for each model. The search was performed using TF-IDF as the
reference vectorisation; the same hyperparameters are applied to all three
representations.

\begin{table}[h!]
\centering
\caption{Hyperparameter search (GridSearchCV, 3 stratified folds, macro F1).}
\label{tab:grid}
\begin{tabular}{L{3cm} L{5.5cm} L{3.5cm} C{1.5cm}}
\toprule
\textbf{Model} & \textbf{Hyperparameters evaluated} & \textbf{Best config.} & \textbf{CV\,F1} \\
\midrule
Log. Regression & $C \in \{0.1,\,1.0,\,10.0\}$ & $C=1.0$ & ${\sim}0.64$ \\
Naïve Bayes & \texttt{var\_smoothing} $\in \{10^{-9},10^{-7},10^{-5}\}$ & $10^{-5}$ & ${\sim}0.47$ \\
$k$-NN & $k \in \{3,\,5,\,9\}$ & $k=5$ & ${\sim}0.36$ \\
Decision Tree & \texttt{max\_depth} $\in \{5,10,20\}$; criterion $\in \{\text{gini, entropy}\}$ & gini, depth\,20 & ${\sim}0.49$ \\
SVM$^\star$ & kernel $\in \{\text{linear, rbf}\}$; $C \in \{1.0,\,5.0\}$ & rbf, $C=1.0$ & ${\sim}0.65$ \\
MLP & layers $\in \{(64),(128,\!64),(128,\!64,\!32)\}$; $\alpha \in \{10^{-4},10^{-3}\}$ & $(128,\!64,\!32)$, $\alpha=10^{-3}$ & ${\sim}0.62$ \\
\bottomrule
\multicolumn{4}{l}{\normalsize $^\star$ Best macro F1 in cross-validation.}
\end{tabular}
\end{table}

\subsection{Test-Set Results --- TF-IDF}

Table~\ref{tab:tfidf} shows performance with TF-IDF (unigrams + bigrams +
trigrams). The CV\,F1 column is the mean of 5-fold cross-validation and is
the primary generalisation indicator; F1\,(train) is included only as a
reference for diagnosing overfitting.

\begin{table}[h!]
\centering
\caption{Results with TF-IDF (1,3) --- test set, macro-averaged metrics.}
\label{tab:tfidf}
\begin{tabular}{L{3cm} C{1.4cm} C{1.4cm} C{1.4cm} C{1.4cm} C{1.6cm} C{1.4cm} C{1.4cm}}
\toprule
\textbf{Model} & \textbf{Acc.} & \textbf{Prec.} & \textbf{Rec.} & \textbf{F1} & \textbf{F1\,(tr.)} & \textbf{CV\,F1} & \textbf{CV\,std} \\
\midrule
Log. Regression & 0.703 & 0.666 & 0.686 & 0.665 & 0.770 & 0.648 & 0.018 \\
Naïve Bayes     & 0.485 & 0.532 & 0.503 & 0.474 & 0.604 & 0.461 & 0.022 \\
$k$-NN          & 0.394 & 0.515 & 0.364 & 0.353 & 0.504 & 0.340 & 0.025 \\
Decision Tree   & 0.552 & 0.524 & 0.511 & 0.511 & 0.694 & 0.495 & 0.020 \\
SVM$^\star$     & \textbf{0.722} & \textbf{0.678} & \textbf{0.670} & \textbf{0.674} & 0.923 & \textbf{0.659} & 0.016 \\
MLP             & 0.704 & 0.661 & 0.671 & 0.663 & 0.750 & 0.645 & 0.019 \\
\bottomrule
\multicolumn{8}{l}{\normalsize $^\star$ Best test F1. All metrics are macro-averaged over five classes.}
\end{tabular}
\end{table}

\subsection{Test-Set Results --- Word2Vec CBOW}

\begin{table}[h!]
\centering
\caption{Results with Word2Vec CBOW (\texttt{sg=0}) --- test set, macro-averaged metrics.}
\label{tab:cbow}
\begin{tabular}{L{3cm} C{1.4cm} C{1.4cm} C{1.4cm} C{1.4cm} C{1.6cm} C{1.4cm} C{1.4cm}}
\toprule
\textbf{Model} & \textbf{Acc.} & \textbf{Prec.} & \textbf{Rec.} & \textbf{F1} & \textbf{F1\,(tr.)} & \textbf{CV\,F1} & \textbf{CV\,std} \\
\midrule
Log. Regression & 0.641 & 0.618 & 0.625 & 0.617 & 0.695 & 0.601 & 0.022 \\
Naïve Bayes     & 0.462 & 0.501 & 0.476 & 0.451 & 0.578 & 0.438 & 0.028 \\
$k$-NN          & 0.381 & 0.490 & 0.352 & 0.340 & 0.487 & 0.325 & 0.030 \\
Decision Tree   & 0.521 & 0.498 & 0.485 & 0.484 & 0.668 & 0.468 & 0.024 \\
SVM$^\star$     & \textbf{0.672} & \textbf{0.645} & \textbf{0.638} & \textbf{0.638} & 0.882 & \textbf{0.621} & 0.019 \\
MLP             & 0.661 & 0.632 & 0.641 & 0.633 & 0.722 & 0.615 & 0.021 \\
\bottomrule
\multicolumn{8}{l}{\normalsize $^\star$ Best test F1 for this vectorisation.}
\end{tabular}
\end{table}

\subsection{Test-Set Results --- Word2Vec Skip-gram}

\begin{table}[h!]
\centering
\caption{Results with Word2Vec Skip-gram (\texttt{sg=1}) --- test set, macro-averaged metrics.}
\label{tab:sg}
\begin{tabular}{L{3cm} C{1.4cm} C{1.4cm} C{1.4cm} C{1.4cm} C{1.6cm} C{1.4cm} C{1.4cm}}
\toprule
\textbf{Model} & \textbf{Acc.} & \textbf{Prec.} & \textbf{Rec.} & \textbf{F1} & \textbf{F1\,(tr.)} & \textbf{CV\,F1} & \textbf{CV\,std} \\
\midrule
Log. Regression & 0.648 & 0.621 & 0.629 & 0.621 & 0.702 & 0.606 & 0.021 \\
Naïve Bayes     & 0.468 & 0.507 & 0.481 & 0.457 & 0.581 & 0.442 & 0.027 \\
$k$-NN          & 0.385 & 0.492 & 0.358 & 0.344 & 0.490 & 0.329 & 0.029 \\
Decision Tree   & 0.528 & 0.503 & 0.491 & 0.489 & 0.671 & 0.472 & 0.023 \\
SVM$^\star$     & \textbf{0.679} & \textbf{0.649} & \textbf{0.643} & \textbf{0.643} & 0.888 & \textbf{0.626} & 0.018 \\
MLP             & 0.666 & 0.635 & 0.644 & 0.637 & 0.728 & 0.619 & 0.020 \\
\bottomrule
\multicolumn{8}{l}{\normalsize $^\star$ Best test F1. Skip-gram consistently outperforms CBOW across all models.}
\end{tabular}
\end{table}

\subsection{F1 Macro Comparison --- All Three Vectorisations}

Table~\ref{tab:comparison} consolidates F1 macro scores across models and
vectorisations. The best overall combination is \textbf{SVM + TF-IDF\,(1,3)}
with F1\,=\,0.674.

\begin{table}[h!]
\centering
\caption{F1 macro comparison across models and vectorisations.}
\label{tab:comparison}
\begin{tabular}{L{3cm} C{2.2cm} C{2.2cm} C{2.2cm} C{2.2cm} C{2.5cm}}
\toprule
\textbf{Model} & \textbf{TF-IDF (1,1)} & \textbf{TF-IDF (1,2)} & \textbf{TF-IDF (1,3)} & \textbf{W2V CBOW} & \textbf{W2V Skip-gram} \\
\midrule
Log. Regression & 0.658 & 0.662 & 0.665 & 0.617 & 0.621 \\
Naïve Bayes     & 0.468 & 0.471 & 0.474 & 0.451 & 0.457 \\
$k$-NN          & 0.345 & 0.349 & 0.353 & 0.340 & 0.344 \\
Decision Tree   & 0.504 & 0.508 & 0.511 & 0.484 & 0.489 \\
SVM$^\star$     & 0.668 & 0.671 & \textbf{0.674} & 0.638 & 0.643 \\
MLP             & 0.657 & 0.660 & 0.663 & 0.633 & 0.637 \\
\bottomrule
\multicolumn{6}{l}{\normalsize $^\star$ Best global combination: SVM + TF-IDF\,(1,3), F1\,=\,0.674.}
\end{tabular}
\end{table}

\subsection{Detailed Report --- Best Model}

The best model is SVM with TF-IDF\,(1,3) (macro F1\,=\,0.674).
Table~\ref{tab:report} breaks down performance by class on the 1,200-review
test set (240 per class).

\begin{table}[h!]
\centering
\caption{Classification report --- SVM + TF-IDF\,(1,3).}
\label{tab:report}
\begin{tabular}{L{3.5cm} C{2.2cm} C{2.2cm} C{2.2cm} C{2.2cm}}
\toprule
\textbf{Class} & \textbf{Precision} & \textbf{Recall} & \textbf{F1} & \textbf{Support} \\
\midrule
very bad  & 0.86 & 0.83 & 0.85 & 240 \\
bad       & 0.55 & 0.50 & 0.52 & 240 \\
neutral   & 0.65 & 0.63 & 0.64 & 240 \\
good      & 0.68 & 0.72 & 0.70 & 240 \\
very good & 0.72 & 0.69 & 0.70 & 240 \\
\midrule
macro avg & 0.69 & 0.67 & 0.67 & 1,200 \\
\bottomrule
\end{tabular}
\end{table}

The \textit{bad} class shows the lowest F1 (0.52) across all vectorisations.
The boundary between mid-range polarity and adjacent extremes (\textit{very
bad} and \textit{neutral}) is inherently ambiguous; the LLM ensemble disagrees
frequently in that region, yielding fewer training examples for that category.

\subsection{Overfitting Analysis and Cross-Validation}

With TF-IDF, the SVM shows the largest gap between training F1 (0.923) and
test F1 (0.674). However, its CV\,F1 of 0.659 is stable across folds
(standard deviation 0.016), indicating that the model generalises well despite
that gap.

When cross-validation is available, CV\,F1 is the primary generalisation
indicator. The train--test gap only carries diagnostic weight when CV\,F1 is
also high \emph{and} unstable across folds. A large gap with a consistent
CV\,F1 signals partial memorisation, not a generalisation failure.

With Word2Vec, models show smaller gaps because the 100-dimensional
representation acts as an implicit regulariser. The exception is the Decision
Tree (gap ${\approx}0.18$ with CBOW), which memorises dense values without
sufficient depth restriction.

% ─────────────────────────────────────────────────────────────────────────────
\section{Conclusions}

SVM with TF-IDF trigrams achieved the highest macro F1 (0.674) on the
1,200-review test set. Logistic Regression, with F1\,=\,0.665 and a gap of
only 0.105, is the more stable option for classifying unseen reviews in
production.

TF-IDF outperforms Word2Vec on this corpus. Short reviews with technical
vocabulary---seed brand names, agronomic terms---are better represented by
frequency-based weights than by dense semantic vectors. The difference is
consistent across all six classifiers.

Skip-gram systematically outperforms CBOW by a small margin (0.003--0.006 F1).
Low-frequency technical words in this domain are better represented by the
Skip-gram architecture, which optimises the representation of rare words.

The \textit{bad} category is the weakest point of the classifier for all
vectorisations (F1 between 0.48 and 0.52). The LLM ensemble rarely agrees on
mid-range polarity examples, leaving fewer training instances for that
category. Expanding the labelled corpus or applying targeted human review of
this boundary region would likely improve performance on this class.

% ─────────────────────────────────────────────────────────────────────────────
\section*{References}
\addcontentsline{toc}{section}{References}

\begin{thebibliography}{9}

\bibitem{spacy}
  M.\ Honnibal and I.\ Montani,
  ``spaCy: Industrial-strength natural language processing in Python,''
  \textit{Software}, \url{https://spacy.io}, 2017.

\bibitem{word2vec}
  T.\ Mikolov, K.\ Chen, G.\ Corrado, and J.\ Dean,
  ``Efficient estimation of word representations in vector space,''
  \textit{arXiv preprint arXiv:1301.3781}, 2013.

\bibitem{sklearn}
  F.\ Pedregosa et al.,
  ``Scikit-learn: Machine learning in Python,''
  \textit{Journal of Machine Learning Research}, vol.~12,
  pp.~2825--2830, 2011.

\bibitem{navarro2023}
  J.\ Navarro Bellido,
  ``Modelos de aprendizaje automático en análisis de sentimiento,''
  Master's thesis, Universitat Politècnica de València, 2023.

\bibitem{lovera2023}
  F.\,A.\ Lovera and Y.\ Cardinale,
  ``Análisis de sentimientos en Twitter: un estudio comparativo,''
  \textit{Revista Científica de Sistemas e Informática}, vol.~3, no.~1,
  e418, 2023.
  DOI: \href{https://doi.org/10.51252/rcsi.v3i1.418}{10.51252/rcsi.v3i1.418}

\bibitem{cortes1995}
  C.\ Cortes and V.\ Vapnik,
  ``Support-vector networks,''
  \textit{Machine Learning}, vol.~20, no.~3, pp.~273--297, 1995.

\end{thebibliography}

\end{document}
